\chapter{The brief, reviewed}\label{ch:brief}

The brief (\file{.claude/prompts/blueprint-spine-review.md}) asked, as its first step, for a
critical reading of itself. This appendix records that reading, the assumptions the review
adopted where the brief left a choice open, what happened during the session that the brief
did not foresee, and what a future run of the same brief would gain from a few changes.

\section{What the brief settles well}

The brief fixes the object (the blueprint and the code built from it), the three criteria in
priority order, the revisions of every source, the rule that the pinned tex is the
specification and the PDF is not to be cited, and the limits of the session (no edit to a
tracked file, nothing posted, no checkout moved). It separates the pin of a dependency from its
upstream \code{main}, and it names the phases, the master theorems and the vocabulary the
report should use. All of this was followed as written.

\section{Choices the brief leaves open, and the choice made}

\begin{description}
\item[Which source is \enquote{deployed}.] The brief treats the specification, the Rust and
  the Python verifier together as the ground truth, asks that a disagreement between them be
  recorded, and asks of every theorem whether it describes \enquote{the deployed protocol or
  a repaired one}. It does not say which source \emph{is} the deployed protocol. The review
  takes the Rust verifier at the pin to be the deployed protocol and the tex to be the
  specification (\cref{sec:vocabulary}). The distinction mattered at once: the public-input
  check differs between them, and the phase built on \code{main} follows the specification.
\item[Which source the Lean should follow when they disagree.] The brief does not say. The
  repository's convention (\enquote{Category A} content is written from the specification and
  then diffed against the code) and the owner's instruction that built the public-input phase
  (\enquote{as the specification writes it}) favour the specification; the review's finding is
  that this produced a theorem about a verifier nobody runs. The report records each such
  disagreement, and recommends the rule of \cref{sec:opt-tradeoff}: model the deployed
  verifier's checks, take the arguments from the specification, report the difference upstream.
\item[VCVio.] The brief lists local checkouts of ArkLib, CompPoly and Clean beside the
  repository; none was listed for VCVio (one exists, at the new pin). Its pinned sources under
  \file{.lake/packages/VCVio/} were read; its upstream \code{main} from a shallow clone made
  under the review's own directory.
\item[The build.] The brief expects probes that build. The main checkout's build was stale
  (it predated Layer~1 and the public-input phase), so the review rebuilt \code{main} once at
  the start (3280 jobs, exit 0). The owner later asked that no further expensive build be run
  on this machine, which is memory limited; when the upgrade of \code{main} then made every
  probe impossible, the owner authorised one more build (3452 jobs, exit 0), and every probe
  was re-run at the new pins.
\item[Where the report lives.] The brief says anywhere convenient, never committed. Mid-session
  the owner asked for the work to be placed on a separate branch and, on being asked, for the
  report to be committed there. The review branch is \code{docs/protocol-blueprint-review},
  created from \code{main} at \code{b435631}; no tracked file of \code{main} is changed on it,
  and it is not pushed.
\item[Secondary repositories.] The comparable verification efforts were consulted through the
  web only (\cref{ch:options}), since nothing in them needed to be built or run.
\item[Models.] The owner asked, after the first wave, that simpler tasks run on a cheaper
  model. The ground truth, the code audits, the boundary, the libraries' adequacy, the
  obligations tree and the hostile re-derivation ran on the review's own model; the citation
  checks, the LaTeX transcriptions, the register, the probe re-runs, the documentation
  inventory and the literature survey ran on Opus.
\end{description}

\section{What happened that the brief did not foresee}

\begin{itemize}
\item A usage limit stopped eleven of twelve agents an hour into the first wave. The three that
  had written their dossiers incrementally lost nothing; the others lost their conclusions and
  were resumed from their transcripts. The rule \enquote{write the dossier incrementally} was
  added to the shared brief and held from then on.
\item The owner upgraded \code{main} to Lean 4.34.1 with new library pins and merged it into
  the review branch while eleven agents were reading. The proof-system statements did not
  change; the library sources under \file{.lake/packages/} did, so every library citation now
  names its revision and old-pin files are read from git; a numeral in $\K$ changed meaning,
  and every probe was re-run.
\item The leanVM checkout moved off the pin during the same upgrade (a pull and a branch
  checkout by the owner). The pin is an ancestor of the new head, so every leanVM file was
  read at the pin through git from then on, and the citation checks re-read the earlier
  dossiers' lines that way.
\item Two agents forked helpers of their own; the owner stopped two of those, and the rule
  \enquote{no sub-agents} was added for the rest.
\end{itemize}

\section{Scope the brief does not name, which the review covered}

\begin{itemize}
\item \textbf{The setup.} The announced sizes, their canonical encoding, the caps and the
  stacking window are checks of the deployed verifier that precede every phase; the brief's
  list of phases omits them, and the check-by-check enumeration of \cref{ch:drift} includes
  them.
\item \textbf{Ring switching}, a step of the Flock phase with its own challenges and error.
\item \textbf{The generic components} (the sumcheck, the batching, the GKR), which three
  phases are built from and where two of the major findings sit.
\item \textbf{The error accounting}, where the review found the spine's error to be the
  phases' own declaration, and \textbf{the shape of the non-interactive theorems}, which the
  brief does not ask about and which the review found not well formed.
\end{itemize}

\section{What a future run of the brief would gain}

\begin{enumerate}
\item Say which source is the deployed protocol and which source the Lean follows when the
  sources disagree, and what is then owed.
\item Say whether probes may copy repository modules; the mutation probes of the public-input
  phase could only be built as copies of the whole module, its lemmas being private.
\item Name the setup and the generic components among the things to review.
\item Give the expected form of the catalogue: a catalogue of every declaration with its
  snippet is several hundred pages, and the review kept full snippets for the load-bearing and
  interface declarations and one line for the rest.
\item State the budget: how many agents may run at once and on which model. The first wave
  exhausted the session's allowance in under an hour.
\item Freeze the checkouts for the session's duration, or say what to do when they move.
\end{enumerate}
